% predictions.tex --- Section 6: Quantitative Predictions
% ASSERT_CONVENTION: natural_units=kB_equals_1, entropy_base=nats, state_normalization=Tr(rho)=1, von_neumann_entropy=S(rho)=-Tr(rho*ln*rho), information=I(B;M)=S(B)+S(M)-S(BM), free_energy=F=E-TS

%----------------------------------------------------------------------
\section{Quantitative Predictions}\label{sec:predictions}
%----------------------------------------------------------------------

The entropy gradient theorem (Sec.~\ref{sec:gradient}) and the
Landauer bound (Sec.~\ref{sec:landauer}) together constrain the
thermodynamic environment in which self-modeling composites can exist.
We now extract quantitative predictions from these results.  The
predictions fall into two classes: \emph{framework predictions}
(P1--P6), which follow from assumptions A1--A7 alone and carry no
model-dependent parameters, and \emph{order-of-magnitude estimates}
(P7--P9), which additionally require cosmological inputs (A8--A10).
We compare the Landauer-derived entropy deficit with Penrose's
gravitational entropy estimate~\cite{Penrose1979}, finding a gap of
94~orders of magnitude that we argue is a feature of the framework
rather than a failure.

%----------------------------------------------------------------------
\subsection{Minimum entropy deficit for self-modeling}
\label{sec:min-entropy}
%----------------------------------------------------------------------

From the exhaustion bound (Corollary~\ref{cor:exhaustion}), a
self-modeler with initial free energy surplus
$\Delta F = F(\rho_0) - F_{\mathrm{eq}}$ can sustain at most
$N_{\mathrm{cycles}} \lesssim \Delta F / (\kB T\,\MI{B}{M})$
cycles.  Since the available free energy is
$\Delta F \approx \kB T\,\Delta S$ with
$\Delta S = \Smax - S$ the entropy deficit (in nats)---an
approximation valid when the energy spectrum is much smaller
than $\kB T$ (the exact expression is
$\Delta F = \kB T\,D(\rho\|\rho_{\mathrm{eq}})$, where $D$
denotes relative entropy)---the temperature cancels:
\begin{equation}
  N_{\mathrm{exhaust}} = \frac{\Delta S}{\MI{B}{M}} \,.
  \label{eq:n-exhaust}
\end{equation}
This result is purely informational: neither the bath temperature~$T$
nor the coupling constant~$J$ appears.  For self-modeling to occur at
all ($N_{\mathrm{exhaust}} \geq 1$), the initial entropy must satisfy
\begin{equation}
  \boxed{S_{\mathrm{initial}} \leq \Smax - N\,\MI{B}{M}}
  \label{eq:s-initial-bound}
\end{equation}
for $N$~update cycles.  The minimum entropy deficit is therefore
\begin{equation}
  \Delta S_{\min} = N\,\MI{B}{M} \,.
  \label{eq:delta-s-min}
\end{equation}
This is an exact consequence of the Landauer bound
(Theorem~\ref{thm:landauer}) and the definition of available work.
It requires only assumptions A1~(finite-dimensional QM) and
A2~(thermal contact).

%----------------------------------------------------------------------
\subsection{Comparison with the Penrose entropy deficit}
\label{sec:penrose}
%----------------------------------------------------------------------

To evaluate Eq.~\eqref{eq:delta-s-min} numerically, we introduce
three cosmological inputs that are \emph{external} to the framework.

\paragraph{A8 (cosmological entropy budget).}
The Bekenstein--Bousso bound on the cosmological horizon gives
\begin{equation}
  \Smax \sim \frac{\pi R_H^2}{\ell_P^2} \sim 10^{122}\;\text{nats}\,,
  \label{eq:smax-cosmo}
\end{equation}
where $R_H \sim 4.4\times 10^{26}\,\text{m}$ is the Hubble radius
and $\ell_P \sim 1.6\times 10^{-35}\,\text{m}$ is the Planck
length.

\paragraph{Penrose initial entropy.}
Penrose estimated the initial gravitational entropy of the observable
universe from the Weyl curvature hypothesis~\cite{Penrose1979}:
\begin{equation}
  S_{\mathrm{initial}} \sim 10^{88}\;\text{nats}\,,
  \label{eq:penrose-s}
\end{equation}
yielding a Penrose entropy deficit
$\Delta S_{\mathrm{Penrose}} = \Smax - S_{\mathrm{initial}}
\approx 10^{122}\;\text{nats}$ (since $10^{122} \gg 10^{88}$).

\paragraph{A9 (neural self-modeling dimension) and A10 (update rate).}
A human-scale self-modeler with $\sim 10^{10}$ neural degrees of
freedom at $\sim\!\ln 2$~nats each has mutual information
$\MI{B}{M} \sim 7\times 10^{9}\;\text{nats}$.  At an update rate
$\nu \sim 10\;\text{Hz}$ over a cosmic time
$t_U \sim 4\times 10^{17}\;\text{s}$, the total number of cycles is
$N \sim \nu\, t_U \sim 4\times 10^{18}$.

\medskip
The Landauer entropy deficit is then
\begin{equation}
  \Delta S_{\mathrm{Landauer}}
  = N\,\MI{B}{M}
  \sim 4\times 10^{18}\times 7\times 10^{9}
  \sim 3\times 10^{28}\;\text{nats}\,.
  \label{eq:delta-s-landauer}
\end{equation}
The ratio to Penrose's deficit is
\begin{equation}
  \frac{\Delta S_{\mathrm{Landauer}}}{\Delta S_{\mathrm{Penrose}}}
  \sim \frac{3\times 10^{28}}{10^{122}} \sim 10^{-94}\,.
  \label{eq:gap-94}
\end{equation}
The Landauer bound on self-modeling requires an entropy deficit of
order $10^{28}$~nats.  The observed deficit (as estimated by
Penrose) is $\sim\!10^{122}$~nats---larger by
\textbf{94~orders of magnitude}.

This gap is a \emph{feature}, not a failure, of the framework.  The
Landauer bound constrains the \emph{information-processing} cost of
self-modeling: how much entropy headroom a self-modeler must consume
to maintain body--model correlations.  Penrose's estimate constrains
the \emph{gravitational} entropy budget: how smooth the initial
geometry must be relative to the black-hole-dominated equilibrium.
These are different physical mechanisms.  The enormous gap reflects
the fact that gravitational entropy dominates the universe's entropy
budget, while the thermodynamic cost of self-modeling is negligible by
comparison.

The gap is robust to order-of-magnitude variations in the
model-dependent parameters.  Even if $\MI{B}{M}$ is uncertain by a
factor of~$10^3$ and $N$~by a factor of~$10^3$, the Landauer deficit
shifts to~$\sim\!10^{34}$~nats, leaving a gap of~$\sim\!88$ orders
of magnitude.

%----------------------------------------------------------------------
\subsection{Exhaustion timescale}\label{sec:exhaust-time}
%----------------------------------------------------------------------

The exhaustion timescale---the time before the entropy budget is
consumed and self-modeling becomes thermodynamically impossible---is
\begin{equation}
  \tauex = N_{\mathrm{exhaust}}\,\tcyc
  = \frac{\Delta S}{\MI{B}{M}}\,\tcyc \,.
  \label{eq:tau-exhaust}
\end{equation}
Using the Penrose deficit $\Delta S \sim 10^{122}$~nats,
$\MI{B}{M} \sim 7\times 10^{9}$~nats, and
$\tcyc = 1/\nu \sim 0.1\;\text{s}$:
\begin{equation}
  \tauex \sim \frac{10^{122}}{7\times 10^{9}}
    \times 0.1\;\text{s}
  \sim 10^{111}\;\text{s}\,.
  \label{eq:tau-value}
\end{equation}
Compared with the current age of the universe,
\begin{equation}
  \frac{\tauex}{t_U} \sim \frac{10^{111}}{4\times 10^{17}}
  \sim 10^{93}\,.
  \label{eq:tau-ratio}
\end{equation}
Self-modeling is thermodynamically permitted for a time
$\sim\!10^{93}$ times the present age of the universe.  The entropy
budget is far from exhausted; the universe is in the earliest stages
of its thermodynamic lifetime as measured by the self-modeling
capacity.

Note that $\tauex$ depends on the entropy budget (A8--A9) and the
self-modeling parameters (A10), but does \emph{not} depend on the
bath temperature~$T$, which cancels between the available work
$\kB T\,\Delta S$ and the cost per cycle $\kB T\,\MI{B}{M}$.

%----------------------------------------------------------------------
\subsection{Experiential density profile}\label{sec:rho-profile}
%----------------------------------------------------------------------

The experiential density~\cite{Paper5},
\begin{equation}
  \rhoexp(\MI{B}{M})
  = \MI{B}{M}\!\left(1 - \frac{\MI{B}{M}}{S_B}\right),
  \label{eq:rho-profile}
\end{equation}
where $S_B = \SvN{\rho_B}$ is the body entropy, is a concave
parabola in $\MI{B}{M}$ on the interval $[0,\, S_B]$.  Setting
$\mathrm{d}\rhoexp/\mathrm{d}I = 1 - 2I/S_B = 0$ yields
\begin{equation}
  I^{*} = \frac{S_B}{2}\,,
  \qquad
  \rhoexp_{\max} = \frac{S_B}{4}\,.
  \label{eq:rho-peak}
\end{equation}

The profile has three notable limits:
\begin{enumerate}
  \item $\MI{B}{M} = 0$ (no tracking):
    $\rhoexp = 0$.  At thermal equilibrium, the model contains no
    information about the body
    (Eq.~\eqref{eq:eq-mutual}).
  \item $\MI{B}{M} = S_B$ (saturated tracking):
    $\rhoexp = 0$.  The pure-state constraint
    $S(\rho_{BM}) = 0$ enforces $\MI{B}{M} = 2\,S_B$ only for
    pure joint states; for $\MI{B}{M} = S_B$ the parabola returns
    to zero.
  \item $N \to \infty$ (equilibration):
    The monotonicity chain~\eqref{eq:monotonicity} drives
    $\MI{B}{M}\to 0$, so $\rhoexp\to 0$ regardless of initial
    conditions.
\end{enumerate}

Under the passive equilibration dynamics of
Sec.~\ref{sec:three-regimes}, the mutual information decays
geometrically: $I(N) \sim I_0\,(1-p)^N$, where
$p = \sin^2(Jt)$ is the depolarizing parameter.
The time dependence of~$\rhoexp$ then takes the form
\begin{equation}
  \rhoexp(N) = I_0\,(1-p)^N
    \left(1 - \frac{I_0\,(1-p)^N}{S_B}\right).
  \label{eq:rho-time}
\end{equation}
If $I_0 > S_B/2$, this rises to the peak $\rhoexp_{\max} = S_B/4$
(when $I$ has decayed to $S_B/2$) and then falls to zero.  If
$I_0 \leq S_B/2$, $\rhoexp$ decreases monotonically from the start.
Both scenarios terminate at $\rhoexp = 0$ as the system
equilibrates.

%----------------------------------------------------------------------
\subsection{Model-dependence register}\label{sec:model-dep}
%----------------------------------------------------------------------

The quantitative predictions P7--P9 rely on three assumptions beyond
the framework (A1--A7).  We record each, its role, and the sensitivity
of predictions to its value.

\begin{center}
\begin{tabular}{@{}llll@{}}
  \toprule
  ID & Parameter & Enters via & Sensitivity \\
  \midrule
  A8 & $\Smax\sim 10^{122}$ nats
     & Entropy budget $\Delta S$
     & $\tauex \propto \Smax$ \\
  A9 & $S_{\mathrm{init}}\sim 10^{88}$ nats
     & Penrose comparison
     & Gap shifts $\pm 2$ orders \\
  A10 & $d_M\!\sim\! 10^{11}$, $\nu\!\sim\! 10\;\text{Hz}$
      & $\MI{B}{M}$, $N$ in Eq.~\eqref{eq:delta-s-landauer}
      & $\Delta S_{\mathrm{L}}\!\propto\! I\!\cdot\! N$ \\
  \bottomrule
\end{tabular}
\end{center}

\noindent
These are \emph{separate} from the framework assumptions A1--A7.
Varying A8--A10 by several orders of magnitude in either direction
shifts the Landauer deficit $\Delta S_{\mathrm{Landauer}}$
accordingly but does not close the 94-order gap with Penrose.  The
framework predictions P1--P6 (Table~\ref{tab:predictions}) are
independent of A8--A10.

%----------------------------------------------------------------------
\subsection{Prediction summary}\label{sec:pred-table}
%----------------------------------------------------------------------

Table~\ref{tab:predictions} collects the predictions of the
entropy gradient theorem and its consequences.  Each entry is
classified by strength and the assumptions it requires.

\begin{table*}[t]
\caption{Predictions of the self-modeling thermodynamics framework.
``Strength'' indicates the epistemic status: THEOREM
(rigorous under stated assumptions), SELECTION (logical chain with
identified weakest link), MATHEMATICAL (follows from definitions),
ESTIMATE (order-of-magnitude, model-dependent), STRUCTURAL
(qualitative connection without a quantitative formula).
Assumptions A1--A7 are framework assumptions
(Sec.~\ref{sec:gradient}); A8--A10 are cosmological inputs
(Sec.~\ref{sec:model-dep}).
\label{tab:predictions}}
\begin{tabular}{@{}clll@{}}
  \toprule
  \# & Prediction & Strength & Assumptions \\
  \midrule
  P1 & Entropy gradient required: $S(t) < \Smax$ for $\rhoexp > 0$
     & THEOREM & A1--A3 \\
  P2 & $\Delta S \geq \MI{B}{M}$ per update cycle
     & THEOREM & A1, A2 \\
  P3 & $\rhoexp = 0$ at thermal equilibrium
     & THEOREM & A1, A5 \\
  P4 & Complexification necessary for SM-like observers (chirality requires $\Spin{10}$);
       non-SM self-modelers in non-complexified blocks remain open
     & THEOREM + OPEN & A1--A7 \\
  P5 & $\rhoexp$ peaks at $I = S_B/2$ with $\rhoexp_{\max} = S_B/4$
     & MATH & A5 \\
  P6 & $\rhoexp \to 0$ at heat death
     & THEOREM & A1, A3, A5 \\
  \midrule
  P7 & $\Delta S_{\mathrm{Landauer}} \sim 10^{28}$ nats
     & ESTIMATE & A1, A2, A8, A10 \\
  P8 & Landauer $\sim 94$ orders weaker than Penrose
     & ESTIMATE & A8--A10 \\
  P9 & $\tauex\sim 10^{111}\;\text{s}\gg t_U$
     & ESTIMATE & A1, A2, A8--A10 \\
  \bottomrule
\end{tabular}
\end{table*}

%----------------------------------------------------------------------
\subsection{Non-claims}\label{sec:non-claims}
%----------------------------------------------------------------------

To guard against overclaiming, we state explicitly what the
quantitative results do \emph{not} establish.

\begin{description}
  \item[NC-1.] The Landauer bound is not a prediction of the initial
    entropy of the universe.  It provides a \emph{necessary
    condition} ($\Delta S \gtrsim 10^{28}$~nats), not a
    determination.  The observed deficit ($\sim\!10^{122}$~nats) is
    vastly larger, and its origin lies in gravitational
    physics~\cite{Penrose1979}, not in the self-modeling framework.

  \item[NC-2.] The comparison with Penrose is not a derivation of
    $S_{\mathrm{initial}} \sim 10^{88}$.  The two constraints
    address different physics: Landauer constrains
    information-processing cost; Penrose constrains gravitational
    entropy via the Weyl curvature hypothesis.  The comparison
    demonstrates internal consistency (the framework's bound is
    weaker), not explanatory power.

  \item[NC-3.] The exhaustion timescale $\tauex \sim 10^{111}\;\text{s}$
    is not a prediction of the age of the universe.  It quantifies how
    long the entropy budget permits self-modeling to continue, not how
    old the universe is.

  \item[NC-4.] The experiential density profile $\rhoexp(I)$ is not
    computed to observational precision.  The qualitative shape (peak
    at $I = S_B/2$, decay to zero at equilibrium) is exact given
    assumption~A5; the mapping from update steps~$N$ to cosmological
    time is model-dependent (requiring A8--A10 and the specific
    equilibration dynamics from Sec.~\ref{sec:entropy}).
\end{description}

\medskip
The quantitative results of this section are best understood as
consistency checks: the framework's thermodynamic constraints are
compatible with---and vastly weaker than---independent cosmological
estimates~\cite{Penrose1979,Albert2000}.  This weakness is itself
informative: the thermodynamic cost of self-modeling is negligible
compared to the gravitational entropy budget, confirming that the
entropy gradient theorem constrains the \emph{qualitative existence}
of an arrow of time, not its quantitative magnitude.
